%%
%% This is file `sample-sigconf.tex',
%% generated with the docstrip utility.
%%
%% The original source files were:
%%
%% samples.dtx  (with options: `all,proceedings,bibtex,sigconf')
%% 
%% IMPORTANT NOTICE:
%% 
%% For the copyright see the source file.
%% 
%% Any modified versions of this file must be renamed
%% with new filenames distinct from sample-sigconf.tex.
%% 
%% For distribution of the original source see the terms
%% for copying and modification in the file samples.dtx.
%% 
%% This generated file may be distributed as long as the
%% original source files, as listed above, are part of the
%% same distribution. (The sources need not necessarily be
%% in the same archive or directory.)
%%
%%
%% Commands for TeXCount
%TC:macro \cite [option:text,text]
%TC:macro \citep [option:text,text]
%TC:macro \citet [option:text,text]
%TC:envir table 0 1
%TC:envir table* 0 1
%TC:envir tabular [ignore] word
%TC:envir displaymath 0 word
%TC:envir math 0 word
%TC:envir comment 0 0
%%
%% The first command in your LaTeX source must be the \documentclass
%% command.
%%
%% For submission and review of your manuscript please change the
%% command to \documentclass[manuscript, screen, review]{acmart}.
%%
%% When submitting camera ready or to TAPS, please change the command
%% to \documentclass[sigconf]{acmart} or whichever template is required
%% for your publication.
%%
%%
\documentclass[sigconf]{acmart}
%%
%% \BibTeX command to typeset BibTeX logo in the docs
\AtBeginDocument{%
  \providecommand\BibTeX{{%
    Bib\TeX}}}

%% Rights management information.  This information is sent to you
%% when you complete the rights form.  These commands have SAMPLE
%% values in them; it is your responsibility as an author to replace
%% the commands and values with those provided to you when you
%% complete the rights form.
\setcopyright{acmlicensed}
\copyrightyear{2026}
\acmYear{2026}
\acmDOI{XXXXXXX.XXXXXXX}
%% These commands are for a PROCEEDINGS abstract or paper.
\acmConference[Conference acronym 'XX]{Make sure to enter the correct
  conference title from your rights confirmation email}{June 03--05,
  2018}{Woodstock, NY}
%%
%%  Uncomment \acmBooktitle if the title of the proceedings is different
%%  from ``Proceedings of ...''!
%%
%%\acmBooktitle{Woodstock '18: ACM Symposium on Neural Gaze Detection,
%%  June 03--05, 2018, Woodstock, NY}
\acmISBN{978-1-4503-XXXX-X/2018/06}


%%
%% Submission ID.
%% Use this when submitting an article to a sponsored event. You'll
%% receive a unique submission ID from the organizers
%% of the event, and this ID should be used as the parameter to this command.
%%\acmSubmissionID{123-A56-BU3}

%%
%% For managing citations, it is recommended to use bibliography
%% files in BibTeX format.
%%
%% You can then either use BibTeX with the ACM-Reference-Format style,
%% or BibLaTeX with the acmnumeric or acmauthoryear sytles, that include
%% support for advanced citation of software artefact from the
%% biblatex-software package, also separately available on CTAN.
%%
%% Look at the sample-*-biblatex.tex files for templates showcasing
%% the biblatex styles.
%%

%%
%% The majority of ACM publications use numbered citations and
%% references.  The command \citestyle{authoryear} switches to the
%% "author year" style.
%%
%% If you are preparing content for an event
%% sponsored by ACM SIGGRAPH, you must use the "author year" style of
%% citations and references.
%% Uncommenting
%% the next command will enable that style.
%%\citestyle{acmauthoryear}


%%
%% end of the preamble, start of the body of the document source.
\begin{document}

%%
%% The "title" command has an optional parameter,
%% allowing the author to define a "short title" to be used in page headers.
\title{Literature Review}

%%
%% The "author" command and its associated commands are used to define
%% the authors and their affiliations.
%% Of note is the shared affiliation of the first two authors, and the
%% "authornote" and "authornotemark" commands
%% used to denote shared contribution to the research.
\author{Pangseen Puarat}
\affiliation{%
 \institution{University of Auckland}
 \city{Auckland}
 \country{New Zealand}}
\email{ppua856@aucklanduni.ac.nz}

\author{Yao Wang}
\affiliation{%
  \institution{University of Auckalnd}
  \city{Auckalnd}
  \country{New Zealand}}
\email{nway918@auckland.ac.nz}

\author{Aabhas Janbandhu}
\affiliation{%
  \institution{University of Auckland}
  \city{Auckland}
  \country{New Zealand}}
\email{ajan119@aucklanduni.ac.}

%%
%% By default, the full list of authors will be used in the page
%% headers. Often, this list is too long, and will overlap
%% other information printed in the page headers. This command allows
%% the author to define a more concise list
%% of authors' names for this purpose.
\renewcommand{\shortauthors}{Trovato et al.}

%%
%% The abstract is a short summary of the work to be presented in the
%% article.
\begin{abstract}
  The rapid advancement of Large Multimodal Models (LMMs) has enabled artificial intelligence to successfully solve complex, 
  image-based computing assessments, such as visual Parsons problems and graph data structure traversals. 
  Historically, visual problems were utilized to mitigate academic integrity issues associated with text-based generation; 
  however, LMMs can now bypass these safeguards using minimal zero-shot prompting. 
  As educators increasingly look to utilize these models as interactive digital tutors, 
  a critical pedagogical requirement emerges: the ability of LMMs to articulate the step-by-step 
  reasoning behind their visual problem-solving processes. This systematic review investigates whether multimodal models 
  can reliably explain their reasoning processes for image-based computing problems. We synthesize recent literature evaluating 
  Multimodal Chain-of-Thought (MCoT) prompting and LMM performance on visual tasks. Our findings indicate that 
  while LMMs demonstrate high problem-solving accuracy, their capacity for generating robust, multi-step explanations 
  often relies on shallow visual reasoning or defaults to textual context. This review outlines the current explanatory 
  capabilities and limitations of LMMs, offering insights into their pedagogical viability and highlighting future directions 
  for multimodal prompt engineering in computing education.
\end{abstract}

%%
%% The code below is generated by the tool at http://dl.acm.org/ccs.cfm.
%% Please copy and paste the code instead of the example below.
%%
\begin{CCSXML}
<ccs2012>
   <concept>
       <concept_id>10003456.10003457.10003527</concept_id>
       <concept_desc>Social and professional topics~Computing education</concept_desc>
       <concept_significance>500</concept_significance>
   </concept>
   <concept>
       <concept_id>10010405.10010489</concept_id>
       <concept_desc>Applied computing~Education</concept_desc>
       <concept_significance>300</concept_significance>
   </concept>
</ccs2012>
\end{CCSXML}

\ccsdesc[1]{Social and professional topics~Computing education}
%%
%% Keywords. The author(s) should pick words that accurately describe
%% the work being presented. Separate the keywords with commas.
\keywords{Generative AI, Large Multimodal Models, Computing Education, 
Prompt Engineering, Multimodal Chain-of-Thought,
 Visual Programming Problems, Academic Integrity}

%%
%% This command processes the author and affiliation and title
%% information and builds the first part of the formatted document.
\maketitle

\section{Introduction}
Recent breakthroughs in generative artificial intelligence have fundamentally transformed computing education. 
Large language models (LLMs) have demonstrated the ability to automatically 
generate source code from natural language descriptions and perform well on typical 
introductory programming problems, such as CS1 exercises~\cite{denny2023conversing}. 
To address the academic integrity concerns associated with these text-to-text models, 
educators have increasingly shifted toward visual-based assessments and diagrams, 
which have historically limited LLMs' problem-solving capabilities~\cite{gutierrez2024seeing}.
This pedagogical strategy is now being challenged by the advent of large multimodal models 
(LMMs) capable of processing both text and images. Recent empirical evaluations demonstrate 
that LMMs can successfully solve complex, image-based computing tasks. For example, LMMs show 
remarkable proficiency in arranging visually diverse Parsons problems without relying on text-based 
snippets~\cite{hou2024more}, and navigating algorithmic traversals on visual graph and 
tree data structures~\cite{gutierrez2024seeing}. Furthermore, new benchmarks reveal that 
LMMs are increasingly being tested on their ability to interpret structured flowcharts~\cite{singh2024flowvqa}, 
and generate code directly from visually rich programming problems containing graphs, maps, and spatial data~\cite{li2024mmcode}.

While the problem-solving accuracy of LMMs on visual tasks is expanding, their pedagogical utility as 
interactive tutoring tools depends heavily on their ability to articulate how a solution was reached. 
In text-based programming, prompt engineering and natural language interactions are crucial for resolving 
errors and fostering computational thinking~\cite{denny2023conversing}. For multimodal tasks, 
eliciting this step-by-step logic is often formalized as Multimodal Chain-of-Thought (MCoT) reasoning, 
which simulates human-like inference by breaking down complex problems into sequential, interpretable steps~\cite{chen2024m3cot}.

However, existing research highlights significant challenges in the explanatory capabilities of LMMs. 
Studies indicate that current MCoT performance often relies on shallow, single-step visual reasoning or 
defaults to textual context rather than genuine multimodal integration~\cite{chen2024m3cot}. 
When faced with complex schematic diagrams~\cite{zhao2025missqa} or flowcharts requiring directional logic~\cite{singh2024flowvqa}, 
models frequently struggle to contextualize visual elements or produce accurate, multi-step explanations. 
Furthermore, the requirement for intense reasoning over text and images simultaneously remains 
a major hurdle for generating executable code from visual inputs. 
Although techniques like weak-supervision-guided step-by-step explanations show promise in 
improving interpretability for simpler image classification tasks~\cite{jiang2025wise}, 
the extent to which these capabilities translate to articulating solutions for complex computing problems is unclear.

Despite rapid advancements in LMMs, a comprehensive understanding of their ability to 
reliably elucidate their problem-solving processes for visual computing tasks remains limited. 
A systematic review of the literature is necessary to evaluate the current state of multimodal reasoning, 
assess the quality of MCoT outputs, and determine the viability of these models as interactive educational tools. 
Therefore, this literature review is guided by the following primary research question:

Can multimodal models explain the step-by-step reasoning behind solutions to image-based computing problems?

\section{Methods}
In accordance with PRISMA standards for systematic reviews, the methodology for 
evaluating the efficacy of Large Multimodal Models (LMMs) in solving 
computing-based visual questions involved a multi-stage selection and synthesis 
process.

\subsection{Paper Selection and Search Strategy}
Primary literature was identified through systematic searches across three major 
digital repositories: the ACM Digital Library, IEEE Xplore, and arXiv~%
\cite{denny2023conversing, gutierrez2024seeing, hou2024more}. 
The search targeted high-impact peer-reviewed research and technical reports 
published between 2023 and 2026 to capture the rapid transition from text-only 
Large Language Models (LLMs) to current frontier Multimodal Foundation Models 
(MFMs) such as GPT-4o, Gemini 1.5 Pro, and Claude 3.5 Sonnet~%
\cite{gutierrez2024seeing, hou2024more, gupta2024polymath}.

\subsection{Inclusion and Exclusion Criteria}
Papers were selected based on Curricular Grounding, specifically targeting 
``Core Tier 1'' computer science topics including fundamental data structures 
(graphs, trees), algorithms (flowcharts), and visual code logic 
(Parsons problems)~\cite{gutierrez2024seeing, hou2024more}.
\begin{itemize}
    \item \textbf{Inclusion:} Priority was given to research introducing novel, 
    computationally generated benchmarks (e.g., FlowVQA, MISS-QA, POLYMATH) 
    to mitigate ``data leakage'' from existing online educational resources~%
    \cite{singh2024flowvqa, gupta2024polymath, zhao2025missqa}.
    \item \textbf{Exclusion:} Studies focusing exclusively on natural scene 
    photography or text-to-text programming tasks without a visual component 
    were omitted~\cite{gutierrez2024seeing, li2024mmcode}.
\end{itemize}

\subsection{Evaluation Dimensions}
Synthesized studies were evaluated across four critical dimensions for 
diagram-based questions:
\begin{itemize}
    \item \textbf{Structural Complexity:} The impact of variations in node count, 
    edge density, and spatial layout on reasoning accuracy~%
    \cite{singh2024flowvqa, gutierrez2024seeing}.
    \item \textbf{Reasoning Depth:} The distinction between single-hop Fact 
    Retrieval (localized information) and macroscopic Topological Reasoning 
    (structural logic)~\cite{singh2024flowvqa, chen2024m3cot}.
    \item \textbf{Aesthetic Robustness:} Sensitivity to non-structural variations 
    such as node colour, rendered edge width, and font styles~%
    \cite{gutierrez2024seeing, hou2024more}.
    \item \textbf{Visual Grounding (VG):} The ability of models to correctly 
    associate textual generation with specific image regions~%
    \cite{ramachandran2025how, lu2025rsvp}.
\end{itemize}

\section{Results: Objective Synthesis of Multimodal Reasoning Frameworks}
The synthesis of the included studies demonstrates a notable gap in the 
performance levels based on the structural nature and linearity of visual 
computing tasks, in addition to a paradigm shift in the field towards 
high-performing multimodal reasoning frameworks~%
\cite{singh2024flowvqa, gutierrez2024seeing, chen2024m3cot}. 
As the application of MLLMs continues to grow in high-stakes settings, 
the research has shifted from the accuracy of predictions to the transparency 
of reasoning, or whether the model can provide step-by-step explanations 
for diagram-based problem-solving, as verifiable by humans~%
\cite{chen2024m3cot, wang2026bridging, jiang2025wise}.

\subsection{Comparative Analysis of Core Methodologies}
The current research landscape is defined by two complementary directions: 
(i) diagram-based reasoning performance in computing tasks, and 
(ii) interpretability-driven multimodal reasoning frameworks.
On one hand, models such as GPT-4o, Gemini, and Claude are evaluated 
on structured computing diagrams including graphs, trees, and flowcharts~%
\cite{singh2024flowvqa, gutierrez2024seeing}. On the other hand, 
frameworks such as WISE, P2X, and MultiFIX focus on improving reasoning 
transparency through structured decomposition and symbolic alignment~%
\cite{jiang2025wise}.

\begin{table*}
\caption{Comparison of Core Methodologies and Frameworks}
\begin{tabular}{p{0.15\linewidth} p{0.3\linewidth} p{0.25\linewidth} p{0.2\linewidth}}
\toprule
\textbf{Framework} & \textbf{Primary Mechanism} & \textbf{Backbone Models Used} & \textbf{Output Type} \\
\midrule
WISE~\cite{jiang2025wise} & Weak Supervision \& CBM Reformulation & CLIP (ViT-L/14), Qwen2-VL-7B & Multimodal Chain-of-Thought (MCoT) \\
P2X & Program Deconstruction \& Reasoning Cells & VLMO, ALBEF & Bounding Boxes, Classes, Answers \\
MultiFIX & Multimodal Data Fusion + Symbolic Reasoning & DL + Genetic Programming & Symbolic Expressions, Heatmaps \\
\bottomrule
\end{tabular}
\end{table*}

The WISE Logic Hierarchy proposes the following structure in the reasoning pipeline:
\begin{itemize}
    \item \textbf{Prior Tree:} Category-level priors for efficient paths in the reasoning process
    \item \textbf{Affirmation Tree:} Supporting evidence specific to the instance
    \item \textbf{Elimination Tree:} Utilizing counterfactuals to eliminate incorrect paths
\end{itemize}
This hierarchy has close ties to the diagram reasoning task because it involves 
both identifying components and relationships~\cite{jiang2025wise}.

\subsection{Performance on Foundational Computing Diagrams}
MFMs have high skills in solving linear and hierarchical diagrams but face 
difficulties in non-linear diagrams~\cite{gutierrez2024seeing, hou2024more}. 
GPT-4V scored a near-perfect percentage of 96.7\% on visual Parsons problems, 
including all representations such as Runestone and Epplets~\cite{hou2024more}.
In hierarchical data structures, GPT-4o scored the highest at 87.6\% on 
tree-based structures, while there was a drastic fall in accuracy to 44.7\% 
on non-hierarchical graphs, showing a limitation in dealing with complex 
relational dependencies~\cite{gutierrez2024seeing}.

\begin{table*}
\caption{Comparative Model Accuracy (\%) Across Diagrammatic Benchmarks}
\label{tab:accuracy}
\begin{tabular}{lccccc}
\toprule
\textbf{Model Family} & \textbf{Parsons~\cite{hou2024more}} & \textbf{Trees (BST)~\cite{gutierrez2024seeing}} & \textbf{Graphs (UG/DG)~\cite{gutierrez2024seeing}} & \textbf{FlowVQA~\cite{singh2024flowvqa}} & \textbf{Schematics~\cite{zhao2025missqa}} \\
\midrule
GPT-4o/V & 96.7 & 92.5 & 44.7 & 61.2 & 63.0 \\
Gemini 1.5 & 69.2 (Bard) & 70.8 (Pro) & 56.2 (Flash) & 49.5 (Pro) & 67.3 (Flash) \\
Claude 3.5/Opus & --- & 63.1 (Opus) & 25.7 (Opus) & --- & 78.3 (o4-mini) \\
Human Expert & 100.0 & 95.1 & 75.0 & 85.3 & 89.0 \\
\bottomrule
\end{tabular}
\end{table*}

\subsection{Strengths and Failure Modes}
A common trend observed with all diagram-based benchmarks is the existence 
of a topological reasoning deficiency. This is evident with the FlowVQA model, 
where performance is significantly higher on Fact Retrieval (T1) than 
Topological reasoning (T4) tasks. This indicates a deficiency in multi-step 
topological reasoning capabilities. Accuracy is observed to dip drastically 
with an increase in the complexity of the diagrams, reducing from 51\% to 34\% 
with an increase in complexity from simple flowcharts to complex ones~%
\cite{singh2024flowvqa}.

Directional biases point to the deficiency in true visual grounding. 
This is evident with the performance drop of up to 15\% with an inversion 
of the flowcharts from top-down to bottom-up. This indicates that the model 
is not performing true visual grounding but is instead relying on prior 
knowledge about the orientation of the diagrams~%
\cite{singh2024flowvqa, ramachandran2025how}.

An analysis of the errors with the POLYMATH model indicates that around 
60\% of the errors are due to logical errors, while around 25\% are due 
to spatial errors. This indicates a deficiency in reasoning consistency 
and spatial understanding~\cite{gupta2024polymath}.

\subsection{Quantitative Benchmarking of Accuracy and Interpretability}
Experimental results on multimodal reasoning benchmarks 
(CUB, SkinCon, LAD, GQA) show that the structured reasoning framework 
improves both accuracy and interpretability~\cite{jiang2025wise}.

\begin{table}[h]
\caption{Quantitative Benchmarking of Accuracy and Interpretability}
\label{tab:interpretability_benchmarks}
\centering
\begin{tabular}{lccc}
\toprule
\textbf{Dataset} & \textbf{Baseline} & \textbf{MCoT-Guided} & \textbf{Interpretability} \\
\midrule
CUB (Birds) & 82.40\% & 83.69\% & 65.03\% \\
SkinCon-9 & 60.37\% & 62.23\% & 93.45\% \\
LAD-H & 59.02\% & 61.61\% & 83.44\% \\
GQA & 60.78\% & 64.13\% & 85.49\% \\
\bottomrule
\end{tabular}
\end{table}

Key results include:
\begin{itemize}
    \item \textbf{Interpretability improvement:} WISE shows a $\sim$37\% 
    improvement in interpretability compared to zero-shot MCoT~\cite{jiang2025wise}.
    \item \textbf{OOD robustness:} P2X achieves state-of-the-art results 
    on GQA-OOD with up to 4.5\% improvement.
    \item \textbf{Concept efficiency:} $\sim$7 out of 312 concepts are used 
    to classify 200 categories.
\end{itemize}
These results show that reasoning structure contributes to interpretability 
improvement as well as performance improvement.

\subsection{Reasoning Quality and Prompting Strategies}
Prompt engineering continues to play a vital role in improving multimodal 
reasoning efficiency~\cite{singh2024flowvqa, denny2023conversing}. 
Techniques such as Multimodal Chain-of-Thought (M3CoT) can be used to 
decompose complex problems into simpler ones for better reasoning 
efficiency but at a computational cost~\cite{chen2024m3cot}. 
Techniques such as Explainable CoT Compression (XMCC) can be used to 
improve efficiency by reducing redundant steps while still maintaining 
alignment signals~\cite{wang2026bridging}. RSVP can be used to improve 
efficiency by enhancing object granularity for visual tasks~\cite{lu2025rsvp}. 

Reasoning quality can be measured using various metrics beyond accuracy. 
Techniques such as WISE use Gini Impurity to measure decision efficiency, 
while P2X uses BLEU-4, CIDEr-D, and IoU for both linguistic and spatial 
alignment~\cite{jiang2025wise}. 

Ablation studies have shown that the multimodal reasoning pipeline is 
highly sensitive. Removing components such as Salient Concept Selection 
can result in a reduction in interpretability of more than 50\%, while 
removing program disentangling affects accuracy~\cite{jiang2025wise}.

%%
%% The next two lines define the bibliography style to be used, and
%% the bibliography file.
\bibliographystyle{ACM-Reference-Format}
\bibliography{lit-review.bib}


%%
%% If your work has an appendix, this is the place to put it.
\appendix

\end{document}
\endinput
%%
%% End of file `sample-sigconf.tex'.
